% The book's numbers for the 37 labels the guide cites, from the
% 52nd edition's .aux (29 September 2026). Generated by scripts/book_labels.py;
% build.sh appends these lines to the guide's .aux. Do not edit.
\newlabel{def:finite-cell-pattern}{{1.11}{100}{Finite cell pattern}{theorem.1.11}{}}
\newlabel{def:finite-cell-pattern@cref}{{[definition][11][1]1.11}{[1][100][]100}}
\newlabel{def:intwit-book}{{1.9}{88}{Interpretation witness}{theorem.1.9}{}}
\newlabel{def:intwit-book@cref}{{[definition][9][1]1.9}{[1][88][]88}}
\newlabel{def:time-cyclarchy}{{1.8}{62}{Finite cyclarchy}{theorem.1.8}{}}
\newlabel{def:time-cyclarchy@cref}{{[definition][8][1]1.8}{[1][61][]62}}
\newlabel{def:time-history}{{1.6}{57}{Occurrence history}{theorem.1.6}{}}
\newlabel{def:time-history@cref}{{[definition][6][1]1.6}{[1][57][]57}}
\newlabel{def:time-rectangular}{{1.5}{52}{Rectangular tapestry}{theorem.1.5}{}}
\newlabel{def:time-rectangular@cref}{{[definition][5][1]1.5}{[1][52][]52}}
\newlabel{def:time-ticks}{{1.7}{58}{Tick realization}{theorem.1.7}{}}
\newlabel{def:time-ticks@cref}{{[definition][7][1]1.7}{[1][58][]58}}
\newlabel{fig:add-alignment}{{1.9}{14}{The one perfect alignment between $\glue (R_2,R_3)$ and $R_5$. The dashed arrows are the witness: each pairs a tagged cell with a position of the plain rod, and nothing else happens. The upper rod keeps its tags and its seam; read down the arrows and the seam lands between positions $1$ and $2$}{figure.1.9}{}}
\newlabel{fig:add-alignment@cref}{{[figure][9][1]1.9}{[1][14][]14}}
\newlabel{fig:chess-typed-arrows}{{1.57}{126}{Two witnessed legal turns. The labels $\state _W$ and $\state _B$ classify the displayed states; the arrows show these particular legal moves, not total maps on every state of those types. Nothing about the pictures forbids following the first turn with another White turn; the types do. The board is the same board in all three panels --- only the state changes, and with the second turn Black's castling right goes with it, though no panel here shows that either}{figure.1.57}{}}
\newlabel{fig:chess-typed-arrows@cref}{{[figure][57][1]1.57}{[1][125][]126}}
\newlabel{fig:figures-on-board}{{1.48}{108}{Figures on the board. Left: a domino stands on $g7$ and $h7$, two faces that share an edge, and a knight's step shape stands at $d4$ as the eight faces it reaches. Right: the same domino offered at $d6$ and $e7$, which share only a corner. No map sends its link to a shared edge, so the placement type is empty and no witness can be supplied. A figure adds no cell; the map from its cells to the board's is the whole datum}{figure.1.48}{}}
\newlabel{fig:figures-on-board@cref}{{[figure][48][1]1.48}{[1][108][]108}}
\newlabel{fig:form-history-plate}{{1.21}{39}{One tree, two records. The history is the tree itself, atoms at the leaves and constructors at the nodes. The code is what the tree adds up to, and the grey sums read it off node by node. A reduct that keeps the code alone keeps three cells in order and loses every constructor, atom and seam}{figure.1.21}{}}
\newlabel{fig:form-history-plate@cref}{{[figure][21][1]1.21}{[1][39][]39}}
\newlabel{fig:one-code-two-histories}{{1.22}{40}{Two histories over one code. Tree (b) glues two copies of $d$ after $a$ where tree (a) substitutes $d$ into $\glue (b,c)$; both build three cells in order, carrying $a$, $d$ and $d$. The code cannot tell them apart and the history can. The barred arrow marks a pairing that cannot be written: $\History (k')$ holds trees of four cells, and neither of these is one}{figure.1.22}{}}
\newlabel{fig:one-code-two-histories@cref}{{[figure][22][1]1.22}{[1][39][]40}}
\newlabel{fig:time-measurement}{{1.38}{74}{Event association, calibration and elapsed readout in one record. The highlighted columns identify the start and stop samples. Residuals $0$ and $1$ give a pair bound of one reference unit. The source difference yields the displayed interval, which contains the recorded reference difference $3001U$}{figure.1.38}{}}
\newlabel{fig:time-measurement@cref}{{[figure][38][1]1.38}{[1][74][]74}}
\newlabel{fig:time-tapestry}{{1.25}{45}{Four naturally coordinated processes and a clock reading. The palette lists their phases; the coloured rows read them together. The top row stays unchanged, and the others repeat after two, three and four steps. The outlined columns select two joint readings. The separate bottom row names the twelve distinct joint configurations. It records their shared progression without causing their coordination}{figure.1.25}{}}
\newlabel{fig:time-tapestry@cref}{{[figure][25][1]1.25}{[1][44][]45}}
\newlabel{fig:wheel}{{1.15}{25}{Wheel vocabulary. The eight sectors and the cyclic successor are the support; the base cut orders the cycle for reading, so that sector $0$ is where a reading begins; a state selects sectors. Left, the arc state $\{0,1,2\}$; right, the gapped state $\{1,4,6\}$. Both tally $3/8$. The tally is a reading and does not tell them apart; the states do}{figure.1.15}{}}
\newlabel{fig:wheel@cref}{{[figure][15][1]1.15}{[1][25][]25}}
\newlabel{fig:wheel-refinement}{{1.16}{27}{Adding $1/4$ and $1/3$ on one pointed wheel. Refining the four-wheel by $3$ and the three-wheel by $4$ puts both on twelve-sector wheels, with three and four sectors selected from the base cut. Gluing the two counts selects seven sectors from the base cut; the two colours record which summand each sector came from, and the tally $7/12$ is read off the result}{figure.1.16}{}}
\newlabel{fig:wheel-refinement@cref}{{[figure][16][1]1.16}{[1][26][]27}}
\newlabel{prop:rod-laws}{{1.31}{171}{The laws of a rod}{theorem.1.31}{}}
\newlabel{prop:rod-laws@cref}{{[proposition][31][1]1.31}{[1][171][]171}}
\newlabel{prop:tetractys-intwits}{{1.10}{89}{The two tetractys interpretations coexist}{theorem.1.10}{}}
\newlabel{prop:tetractys-intwits@cref}{{[proposition][10][1]1.10}{[1][89][]89}}
\newlabel{sec:add-mult-order}{{1.3}{15}{Addition, multiplication, and order}{section*.17}{}}
\newlabel{sec:add-mult-order@cref}{{[section][3][1]1.3}{[1][15][]15}}
\newlabel{sec:board-canvas}{{1.8}{99}{Constructing a space and its states}{section.1.8}{}}
\newlabel{sec:board-canvas@cref}{{[section][8][1]1.8}{[1][99][]99}}
\newlabel{sec:building-forms}{{1.2}{3}{Building finite forms}{section.1.2}{}}
\newlabel{sec:building-forms@cref}{{[section][2][1]1.2}{[1][3][]3}}
\newlabel{sec:clock-trajectory}{{1.5}{64}{One clock on two scales}{section*.41}{}}
\newlabel{sec:clock-trajectory@cref}{{[section][5][1]1.5}{[1][64][]64}}
\newlabel{sec:complete-assignment}{{1.4}{33}{A complete assignment is a dependent function, and its count multiplies}{section*.28}{}}
\newlabel{sec:complete-assignment@cref}{{[section][4][1]1.4}{[1][33][]33}}
\newlabel{sec:evidence-preview}{{1.4}{42}{Evidence is an input, and its type may be empty}{section*.30}{}}
\newlabel{sec:evidence-preview@cref}{{[section][4][1]1.4}{[1][42][]42}}
\newlabel{sec:formation-lines}{{1.4}{36}{A formation line says what a constructor needs}{section*.29}{}}
\newlabel{sec:formation-lines@cref}{{[section][4][1]1.4}{[1][36][]36}}
\newlabel{sec:glass-bead-game}{{1.7}{93}{The Glass Bead Game}{section.1.7}{}}
\newlabel{sec:glass-bead-game@cref}{{[section][7][1]1.7}{[1][93][]93}}
\newlabel{sec:interpretation}{{1.6}{79}{One artform, many meanings}{section.1.6}{}}
\newlabel{sec:interpretation@cref}{{[section][6][1]1.6}{[1][79][]79}}
\newlabel{sec:rods-grids}{{1.3}{13}{Arithmetic in Artform Calculus}{section.1.3}{}}
\newlabel{sec:rods-grids@cref}{{[section][3][1]1.3}{[1][13][]13}}
\newlabel{sec:role-policy}{{1.9}{121}{Nothing may depend on a reading}{section*.66}{}}
\newlabel{sec:role-policy@cref}{{[section][9][1]1.9}{[1][121][]121}}
\newlabel{sec:rules-counting}{{1.11}{173}{Evidence is counted as a property}{section*.100}{}}
\newlabel{sec:rules-counting@cref}{{[section][11][1]1.11}{[1][172][]173}}
\newlabel{sec:rules-equalities}{{1.11}{172}{Three equalities in \texorpdfstring {$2+3=5$}{2+3=5}}{section*.99}{}}
\newlabel{sec:rules-equalities@cref}{{[section][11][1]1.11}{[1][172][]172}}
\newlabel{sec:rules-rods}{{1.11}{170}{The rules for rods}{section*.98}{}}
\newlabel{sec:rules-rods@cref}{{[section][11][1]1.11}{[1][170][]170}}
\newlabel{sec:rules-view}{{1.11}{168}{Rules of inference: a type-theoretic view of artform \texorpdfstring {\mbox {calculus}~\optionalmark }{calculus}}{section.1.11}{}}
\newlabel{sec:rules-view@cref}{{[section][11][1]1.11}{[1][168][]168}}
\newlabel{sec:shape-type-book}{{1.8}{102}{A square can specify a type}{section*.57}{}}
\newlabel{sec:shape-type-book@cref}{{[section][8][1]1.8}{[1][102][]102}}
\newlabel{sec:tagged-choice}{{1.4}{31}{A tagged choice is a dependent pair, and its count adds}{section*.27}{}}
\newlabel{sec:tagged-choice@cref}{{[section][4][1]1.4}{[1][31][]31}}
\newlabel{sec:time}{{1.5}{43}{Time, clocks, and measurement}{section.1.5}{}}
\newlabel{sec:time@cref}{{[section][5][1]1.5}{[1][43][]43}}
\newlabel{sec:types-values}{{1.4}{30}{A colon and two arrows}{section*.26}{}}
\newlabel{sec:types-values@cref}{{[section][4][1]1.4}{[1][30][]30}}
\newlabel{sec:writing-constructions}{{1.4}{27}{Writing a construction down}{section.1.4}{}}
\newlabel{sec:writing-constructions@cref}{{[section][4][1]1.4}{[1][26][]27}}
